\subsection{Further checks of infinite-ramification and reduction arguments}

\begin{result}{P.12}{Known}{Density-zero reduction is weaker than finite prime support}
For a noncuspidal Selmer section of the \emph{full} arithmetic fundamental
group of a hyperbolic curve over a number field, the primes where its local
point reduces to a cusp have density zero after a suitable finite field
extension. Two such sections with the same reductions on a set of density
one are conjugate. These are
\citet[Theorems~1.1--1.2, pp.1--2; Theorem~5.4 and Corollary~5.7,
pp.14--15]{RootPorowski2025}.
They supply neither a finite-support conclusion nor a theorem for an
arbitrary geometrically pronilpotent section.
\end{result}
\begin{proof}[Derivation and scope]
For the tripod, reduction to its boundary is precisely a nonzero valuation
of one of the two ordinary coordinates, apart from any finitely many
model primes. Thus the stated density theorem concerns the support set
in this paper if a full locally geometric section exists. Its hypotheses
are stronger than ordinary abelian localization of a quotient section.

Even its conclusion is compatible with the constructed ghost families.
In E.102 impose the additional bounds
\(q_j,r_j>\exp(j^2)\). The prime-selection theorem allows any finite
lower bound, so this changes none of that proof. The number of its
support primes at most \(X\) is at most \(2\sqrt{\log X}+O(1)\).
For a fixed number field \(L\), at most \([L:\mathbb Q]\) prime ideals
lie over each rational prime and their norms are at least that prime.
The same estimate, multiplied by \([L:\mathbb Q]\), therefore holds
after base change. Alternatively, the sum of their inverse norms
converges, proving Dirichlet density zero directly. Their support is
still infinite. This is a compatibility check, not a construction of
a full section with that support.
\end{proof}

\begin{result}{P.13}{Known}{An infinite-ramification deformation theorem does not supply the required cusp constraints}
\citet[Theorem~1, p.794; local deformation calculation and
Proposition~1, pp.796--798]{RootRamakrishna2000} constructs infinitely
ramified representations into \(\operatorname{GL}_2(\mathbb Z_p)\)
for a density-one set of primes \(p\ge5\), starting from the stated
semistable elliptic-curve residual representations. The method has
an additional permitted ramified local deformation direction at each
chosen auxiliary prime. It is not a theorem about sections of
\(E_U^{(p)}\), and does not assert an ordinary completed pair.
\end{result}
\begin{proof}[Derivation and scope]
The source fixes a determinant and deforms a surjective residual
two-dimensional representation. Its local ``special'' condition is
\[
 F_q\longmapsto\begin{pmatrix}q&0\\0&1\end{pmatrix},
 \qquad T_q\longmapsto\begin{pmatrix}1&u\\0&1\end{pmatrix}.
\]
The matrix relation \(F_qT_qF_q^{-1}=T_q^q\) holds for every
\(u\in\mathbb Z_p\), which is the available ramified direction.
No comparison homomorphism from these representation spaces to the
tripod's section spaces is supplied. In the present problem, after
fixing a lower integral nilpotent section, corrections in weight
\(w>1\) have coefficients \(\mathbb Z_p(w)\). At an auxiliary prime
\(q\ne p\), their tame inertia
value \(v\) satisfies \((q^{w-1}-1)v=0\); this forces \(v=0\)
because the coefficient module is torsion-free and \(q^{w-1}\ne1\).
Hence importing the displayed auxiliary direction into that fixed
lifting problem would change its coefficient module. A construction
which also changes its lower section is a distinct possibility,
requiring preservation of the old local point constraints at every
step; no such preservation theorem follows from the cited result.
\end{proof}
